\section{Introducción}

El presente capítulo cierra el trabajo contrastando lo construido con lo
comprometido. No se limita a resumir los capítulos precedentes: recorre uno por
uno los objetivos específicos y las hipótesis declarados en el perfil, y para
cada uno establece si la evidencia lo respalda, lo refuta o no alcanza a
pronunciarse. Esa distinción es el aporte central de este capítulo, porque un
sistema complejo rara vez se comporta de forma uniforme: algunos componentes
superan el umbral exigido, otros quedan por debajo y otros no llegaron a
medirse.

El capítulo se organiza en tres partes. La primera sintetiza el sistema
efectivamente construido y delimita qué puede afirmarse de él. La segunda
contrasta la hipótesis general y cada hipótesis específica, y dedica un apartado
propio a la brecha del reconocimiento de matrículas, por ser el componente donde
la distancia entre el criterio comprometido y el resultado medido es mayor. La
tercera formula las recomendaciones que se desprenden de esas mediciones,
distinguiendo las que corresponden a la continuidad inmediata del sistema de
aquellas que constituyen líneas de trabajo posteriores.

Se mantiene a lo largo del capítulo una regla de escritura: ningún resultado se
presenta sin la condición que lo produjo, y ninguna limitación se enuncia sin
indicar qué la causa y qué la removería. Un trabajo que declara sus brechas con
precisión es más útil que uno que las disimula, porque permite que el siguiente
investigador sepa exactamente dónde estaba el límite y por qué.
